Two on-chip memory variants are defined (see Table~\ref{tab:asChipVariants}
and Section~\ref{subsec:chipVariants}):

\begin{itemize}
  \item \textbf{Variant~A (No-Cache):}
    The system memory consists of two on-chip SRAMs.
    The Instruction Memory (32\,kB) is loaded at start-up via the JTAG scan
    chain and is read-only during normal execution.
    The Data Memory (8\,kB) is fully read/write and serves as the sole
    writable data store (stack, heap, \texttt{.data}, \texttt{.bss}).

  \item \textbf{Variant~B (Cache):}
    The system memory consists of an off-chip NOR Flash (code and read-only
    data) plus two on-chip caches and one on-chip Scratchpad SRAM.
    The I-Cache (4\,kB) caches instruction fetches from Flash.
    The D-Cache (4\,kB) caches load accesses to Flash; stores to the Flash
    address region are not permitted and cause a bus error.
    The Scratchpad SRAM (8\,kB) is the sole writable data store and is
    accessed directly (cache bypass) for all load/store operations that
    target its address range.
    The NOR Flash is accessed through the on-chip QSPI controller via the
    AXI4 cache-fill path.
\end{itemize}
